\begin{table}[htbp]
\caption{Comparación de formas funcionales de la edad, por AIC.}
\label{tab:tabla_54_forma_funcional_edad}
\begin{tabular}{llrrrrll}
\toprule
forma & descripcion & aic & gl & log\_verosimilitud & delta\_aic & veredicto & decision \\
\midrule
Spline cúbico & Nudos en [28, 37, 51] años & 14101.845000 & 6 & -7043.922000 & 0.000000 & Mejor AIC & La spline ajusta mejor la edad —72.2 puntos de AIC aislada, 30.6 en el modelo completo— pero no cambia ninguna otra conclusión: el coeficiente que más se mueve lo hace un 4.3 \%, muy por debajo del 10 \% que marca una confusión relevante. Elegir entre una y otra es, por tanto, una decisión de comunicación y no de validez. Se mantiene la categórica: da odds ratios que se leen solos («el grupo de 60+ frente al de 30-39»), es la forma en que reporta la literatura con la que este trabajo se compara, y una spline necesita un gráfico para interpretarse. La spline queda documentada como análisis de sensibilidad. \\
Cúbica & Se permite una curva con dos inflexiones & 14104.646000 & 5 & -7046.323000 & 2.801000 & Peor por 2.8 puntos &  \\
Cuadrática & Se permite una curva en U & 14136.701000 & 4 & -7063.351000 & 34.856000 & Peor por 34.9 puntos &  \\
Categórica (la actual) & Seis grupos, referencia 30-39 & 14174.023000 & 7 & -7079.012000 & 72.178000 & Peor por 72.2 puntos &  \\
Lineal & La edad entra como una recta & 14250.580000 & 3 & -7121.290000 & 148.735000 & Peor por 148.7 puntos &  \\
\bottomrule
\end{tabular}
\end{table}
